\documentclass[../main]{subfiles}
\begin{document}
\chapter{Calidad}
\img{images/09/calidad.png}{Las normas ISO se relacionan con la calidad.}
\section{Herramientas para control de calidad}
\img{images/09/pescado.png}{Diagrama de Ishikawa}
\info{Diagrama de Ishikawa}{El diagrama de Ishikawa, tambión llamado
  de espina de pescado o causa y efecto, es un esquema basado en el
  modelo causal creado en 1943 y perfeccionado por Kaoru Ishikawa. Su
  propósito es organizar el análisis de problemas prioritarios en
  varios procesos, especialmente en producción industrial. Se
  considera una herramienta básica del control de calidad. Este
  diagrama identifica las causas de un problema, agrupándolas en
  categorás para comprender mejor las fuentes de variación.
}
\subsection{Etapas para la elaboración de un diagrama de Ishikawa}
Las principales etapas en la elaboración del diagrama de Ishikawa son
los siguientes:
\begin{enumerate}
  \item Definición del problema.
  \item Elaboración de la representación gráfica.
  \item Análisis de la información que ofrece el diagrama y que
    destaca los principales factores o los factores cuyos valores
    podróan modificarse.
  \item Desarrollar un plan para actuar considerando los comentarios
    de la etapa anterior.
\end{enumerate}
\subsection{Elaboración del diagrama}
Para elaborar el diagrama de Ishikawa, las causas o los factores de
influencia en el problema definido se agrupan en categorás o grupos
principales para identificar estas fuentes de variación y
posteriormente analizarlos. Se construye con mayor frecuencia
utilizando factores o grupos principales de factores de influencia.
En función al tipo de proceso y sector productivo, se utilizan
diferentes métodos que agrupan los factores en cuatro, cinco, seis
o siete componentes, es asóC que sale los métodos de 4M, 5M, 6M o /M.
El más utilizado es el Método de las 6M: que son: mano de obra,
maquinara, materiales, método, medición y medio ambiente.
\subsubsection{Método de las 6M}
Las 6 categorás principales o grupos de influencia se explican de la
siguiente manera:
\begin{itemize}
  \item Mano de obra: Se refiere al talento humano, un factor
    necesario y fundamental en los procesos. El personal puede ser
    operativo, técnico o de gestión, entre otras posibilidades. La
    calidad de la labor de la mano de obra puede estar afectada por el
    nivel de entrenamiento o capacitación, la experiencia del
    personal, la posesión de las competencias requeridas para la
    actividad que realizar, la actitud y el nivel de cansancio.

  \item Maquinaria: Todas las máquinas, equipamiento y tecnologás
    necesarias para realizar las tareas, incluidas las herramientas.
    Incluye sus insumos y actividades de mantenimiento.

  \item Materiales: Se refiere a factores como la materia prima, los
    consumibles o la información requeridos en los procesos.

  \item Método: Se refiere a cómo se lleva a cabo el proceso y los
    requisitos especóficos para hacerlo, como polóticas,
    procedimientos de calidad, instrucciones u órdenes de trabajo,
    planos, normas, reglamentos y leyes.

  \item Medio ambiente: Se enfoca en el anólisis del entorno de
    trabajo de todas las influencias ambientales que pueden ser
    controlables como también las que resultan imprevisibles, como
    ejemplo: tenemos al clima, el tiempo, la temperatura, la humedad,
    las vibraciones, la iluminación, la calidad del aire, el ruido o
    la limpieza que afectan el proceso.

  \item Medición: Es el control, comprobación, evaluación y otras
    medidas fósicas para lograr el proceso, ya sea de manera manual o
    automáticas. Es muy importante que se está atento ante cualquier
    error de calibración u otros problemas de medición para de esta
    manera evitar inconvenientes.
\end{itemize}
Es importante resaltar que, originalmente, se proponen 6 categorás
por este método, sin embargo, no todos los procesos o problemas
utilizan de todos estos factores, asóC que es necesario evaluar cuáles
de ellos están presentes o son importantes para la ejecución. Por
lo cual, es posible que generalmente sólo se evalñen 4 de ellos o, en
 algunos casos, hasta 7.

\subsection{Beneficios}
Como beneficios del uso del diagrama de Ishikawa se pueden mencionar
los siguientes:
\begin{enumerate}
  \item brinda una mayor comprensión de un proceso
  \item permite visualizar información compleja, al organizar y
    relacionar factores, proporcionando una vista secuencial
  \item ayuda a identificar una diversidad de causas a los problemas
    y sus potenciales soluciones;
  \item es un punto de partida para iniciar un proceso de innovación;
  \item facilita un proceso participativo de aprendizaje y el
    intercambio de ideas;
  \item fortalece los vñculos e identidad del equipo de trabajo al
    ser un método participativo;
  \item mejora el manejo de los factores capaces de generar efectos
    menos convenientes;
  \item Identifica las áreas en donde se requiere recoger mayor información;1
  \item identifica la necesidad de elaboración de normas técnicas.
\end{enumerate}

\subsection{Desventajas}

No obstante, el diagrama de Ishikawa puede presentar algunas
desventajas, tales como:

\begin{enumerate}
  \item Dificultad para relacionar algunas causas con una categoráa o
    factor de influencia particular;
  \item El sesgo potencial (cultura, motivación, personalidad y el
    equipo) puede afectar el resultado;
  \item El factor 'mano de obra' suele mal interpretarse: puede
    existir miedo a ser atacado personalmente o  acusado por el
    problema, generando que se genere un barrera para un análisis integral;
  \item No se incluye el tiempo y la evolución del sistema en el
    diagrama, solo presenta un momento especófico de una realidad dinámica;
  \item La necesidad de incluir más factores para adaptar el diagrama al efecto;
  \item Es posible que no se proponga una causa fundamental para
    explicar el efecto.
\end{enumerate}

\actividad{ Investigar y contestar:}
\begin{enumerate}
  \item \textbf{Concepto:} ¿Qué es el Diagrama de Ishikawa y cuál es
    su otro nombre común?
  \item \textbf{Origen:} ¿Quién fue el creador del diagrama de
    Ishikawa y en quá la fue creado?
  \item \textbf{Propósito:} ¿Cuál es la finalidad principal del
    diagrama de Ishikawa en relación con la resolución de problemas?
  \item \textbf{Aplicaciones:} Además de los procesos industriales,
    ¿en qué otros contextos se puede aplicar el diagrama de Ishikawa?
  \item \textbf{Herramientas básicas:} ¿Qué lugar ocupa el diagrama
    de Ishikawa entre las siete herramientas básicas del control de calidad?
  \item \textbf{Definición del problema:} ¿Cuál es la primera etapa
    en la elaboración de un diagrama de Ishikawa y qué se busca
    lograr en esta etapa?
  \item \textbf{Elaboración gráfica:} ¿Qué se hace en la segunda
    etapa para representar las posibles causas de un problema?
  \item \textbf{Análisis de la información:} ¿Qué se destaca en la
    tercera etapa del proceso y qué factores se consideran para posible
    modificación?
  \item \textbf{Plan de acción:} ¿Qué se desarrolla en la cuarta
    etapa y qué se considera para llevar a cabo esta acción?
  \item \textbf{Agrupación de causas:} ¿Cómo se agrupan las causas o
    factores de influencia en el diagrama de Ishikawa?
  \item \textbf{Método de las 6M:} ¿Cuáles son las seis categorás
    principales en el Método de las 6M y qué aspectos abarca cada una?
  \item \textbf{Beneficios:} Menciona al menos tres beneficios del
    uso del diagrama de Ishikawa.
  \item \textbf{Desventajas:} Nombra al menos tres desventajas o
    limitaciones del diagrama de Ishikawa.
\end{enumerate}

\end{document}
